\documentclass[../../main]{subfiles}

\begin{document}

\title{A One-Off Moderate Wealth Tax}
\author{Adrien Fabre}
\date{July 8, 2026}
\maketitle

% The European Union faces a structural financing problem. From 2028 it must begin repaying the common NextGenerationEU (NGEU) debt contracted during the Covid crisis, at the same time as it needs to finance the ecological transition, the competitiveness agenda and a security build-up. The Council remains deadlocked on which new own resources could credibly and acceptably fill the gap: the candidates on the table tend either to raise little, to fall on ordinary households, or to duplicate taxes Member States already levy. This appendix proposes a new own resource that avoids these traps: an exceptional, one-off tax on large individual wealth.


This appendix argues that a new own resource based on wealth taxation can work if designed well, and suggests a way to design it so. %This appendix argues that the EU should collect an exceptional tax on large individual wealth. 
% A wealth tax is an attractive candidate. Large fortunes are today taxed at effective rates well below those borne by ordinary households \citep{Zucman2024}; roughly two thirds of European citizens support taxing the very wealthy \citep{TaxObservatory2025}; and the base, unlike consumption or labour, is concentrated at the very top and, once the levy is designed correctly, barely mobile. %The difficulty is that recurrent wealth taxes have a poor track record, having historically raised little and provoked avoidance and migration. The proposal is built from the outset to avoid these pitfalls.
The proposal deliberately avoids the pitfalls that have historically undermined recurrent wealth taxes.  It is a \emph{one-off} levy, assessed on the net wealth each resident holds on a single past reference date (1~January~2027), and paid over thirty years at low annual marginal rates: 0.1\% on wealth between EUR~1 and 10~million, 0.3\% between EUR~10 and 100~million, and 0.5\% above EUR~100~million. Because the base is fixed in the past, the levy is a lump sum in the economic sense: no decision taken after 2026, whether to invest, to found a company or to emigrate, can change what is owed, so the tax does not distort future accumulation and activity. The low annual rates keep the tax non-confiscatory and comfortably within constitutional bounds. Liquidity is addressed head-on: payment is reduced if a taxpayer's wealth drops and the taxpayer may pay in company shares. Finally, a Deferral Account secures the claim against emigration and death. 

Built this way, the tax would raise %real money. The central scenario yields 
about 
EUR~26~billion a year, close to EUR~790~billion (that is 4.5\% of EU GDP) over the thirty-year horizon. %, in the range of the 5\%-of-GDP benchmark independently proposed by \citet{PisaniFerryMahfouz2023}. % and about twice the common NGEU debt. 
The proposal adopts the low-rate one-off logic of \citet{PisaniFerryMahfouz2023}, who recommended an exceptional levy of 5\% of GDP spread over thirty years to finance the climate transition. 
Besides, it %'s lineage is explicit. It 
adapts the deferral and one-off mechanisms designed for the Californian wealth-tax bill \citep{AB259}. %Unlike %the recurrent net-wealth taxes still levied in Spain, Norway and Switzerland, and unlike 

The instrument rests on solid economic foundations. \citet{Farhi2010} shows that a one-off capital levy is
% in the canonical model of optimal taxation with aggregate shocks and incomplete markets, 
the efficient response to a negative fiscal shock, % is a one-off capital levy, exactly the contingency created by the shocks that motivate new own resources: 
precisely the case created by 
the pandemic-era debt, the security build-up and the financing of the transition. Progressive wealth taxation is independently justified on distributive grounds: \citet{SaezZucman2019} argue that %, by reaching the \emph{stock} of fortune rather than the \emph{flow} of realised income, 
it is the most appropriate tool for restoring progressivity at the very top, since the ultra-wealthy report little taxable income.

% We then argue that, so designed, the tax can actually fly: it does not interact with the taxes Member States already levy; it raises real revenue for national budgets; it is constitutional in every EU jurisdiction; it rests on solid economic and democratic foundations; and it is increasingly on the political agenda. We answer the main objections in turn.

The tax revenue could simply be paid into the EU budget as an own resource, or used to repay the common debt, and we show that it works well on those terms. But we also argue that there is an even better, more strategic use. Rather than shrink the Union's balance sheet by repaying the debt, the proceeds could endow a European sovereign fund, at a time when major investments are needed for the ecological transition and for competitiveness. %Sized at 5\% of EU GDP, the levy would build up a fund whose assets would eventually reach about twice the common debt, so that its long-run dividends would exceed the interest owed several times over. Worldwide, sovereign wealth funds manage assets worth close to 13\% of global GDP. Citizens of some 40 jurisdictions already hold a collective stake in large firms and long-term investments Such funds are no longer the preserve of resource-rich states: China, Singapore, Australia and South Korea each run one worth more than 10\% of their GDP. 
% Yet the European Union has no comparable fund. We develop this option in Section~\ref{sec:fund}.

The combination of low annual rates, deferral, payment in kind and asset-backed recovery is what makes the tax simultaneously enforceable, non-distortive and respectful of taxpayers' liquidity constraints and asset control.
We believe such a tax can actually be adopted. It does not interact with the taxes Member States already levy; it raises real revenue; %for national budgets; 
it is constitutional in every EU jurisdiction; it rests on solid economic and democratic foundations; and it is increasingly on the political agenda.

The remainder covers the proposal (Section~\ref{sec:proposal}), revenue estimates (\ref{sec:revenue}), the case that the tax can fly (\ref{sec:feasibility}), responses to objections (\ref{sec:objections}), the sovereign-fund option (\ref{sec:fund}) and variants (\ref{sec:variants}).

\section{Proposal}
\label{sec:proposal}

\subsection{A one-off tax, payable over thirty years}

We propose a one-off tax on individual net wealth, assessed on the wealth each resident holds on 1 January 2027. The tax is strictly individual: wealth is attributed to natural persons, with couples assessed on their respective shares, so as to avoid the marriage penalties and avoidance opportunities associated with household-level assessment.

The headline schedule is a one-off liability levied at the following marginal rates: 3\% on the fraction of individual wealth between EUR~1 and 10~million, 9\% on the fraction between EUR~10 and 100~million, and 15\% on wealth above EUR~100~million. %The rates apply band by band, so that---as under an income-tax schedule---a taxpayer worth EUR~20~million owes 3\% on the first taxable EUR~9~million and 9\% only on the next EUR~10~million, not 3\% plus 9\% on the whole.
This liability is not due immediately. It is spread over thirty years, so that in any given year the taxpayer owes one thirtieth of the total: the annual tax schedule is 0.1\% on wealth between EUR~1 and 10~million, 0.3\% between EUR~10 and 100~million, and 0.5\% above EUR~100~million. These annual rates are deliberately calibrated well below the return on capital, so that the tax can in principle be paid out of the income the wealth generates, leaving the principal intact.

The one-off character is the central design choice. A permanent wealth tax changes the incentive to accumulate, to invest and to remain resident on an ongoing basis; the empirical literature finds that such taxes can generate %sizeable
behavioural responses, from portfolio reshuffling to migration \citep{Brulhartetal2022, DuranCabreetal2019}. A one-off tax assessed on a past reference date is, by construction, a lump-sum tax: it does not distort decisions taken after the reference date, because no future choice can change the liability. This is the classical efficiency argument for capital levies, and it is the reason \citet{PisaniFerryMahfouz2023} and the Californian bill favoured an exceptional rather than a recurrent instrument.

\subsection{Payment, deferral and the annual minimum}

Each year, the taxpayer must pay the \emph{smaller} of two amounts: (i) one thirtieth of their initial liability, indexed to inflation and increased by a deferral charge; or (ii) the annual schedule (0.1\% / 0.3\% / 0.5\%) applied to their \emph{current} wealth. The second option protects taxpayers whose wealth has fallen: someone who has become poorer since 2027 pays less, in line with their ability to pay. The tax ceases to be owed after thirty years of payment, or earlier if the taxpayer settles their remaining liability in full.

The deferral charge is set at a level that makes postponing payment equivalent for the taxpayer to paying upfront.  
Hence, the deferral charge reflects %to borrowing from the State at a rate that should not be below
the return the taxpayer earns on their wealth. 
We set it at 3\% per year for the tax due for wealth below EUR~10 million, 5\% between EUR~10 million and 100 million, and 7\% above EUR~100 million. These figures are aligned with the empirically documented increasing returns to wealth: larger portfolios earn systematically higher average returns %, a regularity documented across countries and central to the dynamics of wealth concentration
\citep{WorldInequalityReport2022}. % Cite Table 4.1 of WorldInequalityReport2022
Aligning the charge with the return on capital ensures that deferral is neither confiscatory nor a subsidy to the wealthy, and removes the incentive to game the timing of payments.\footnote{Take the example of a taxpayer with a wealth of EUR~10 million in 2027, and assume they have 3\% returns on their wealth and that there is no inflation. The taxpayer's initial tax liability is $3\% \cdot (10-1) = 0.27$M\euro{}. In 2028, if the taxpayer has only paid their annual tax minimum in 2027 (one thirtieth of 270k\euro{}, that is 9,000\euro{}) and not the full tax liability, the taxpayer's expected wealth is $(10,000,000 - 9,000) \cdot(1+3\%)=10,290,730$\euro{}. Their 2028 tax liability is their initial tax liability inflated by the deferral charge minus what  they already paid, that is $(270,000-9,000) \cdot (1+3\%) = 268,830$k\euro{}. Their 2028 wealth minus their remaining liability ($10,290,730-268,830=10,021,900$\euro{}) exactly corresponds to the expected wealth following payment of the full liability in 2027: $(10,000,000-270,000)\cdot (1+3\%)$. Therefore, with a real return of 3\%, the taxpayer is indifferent between paying the entire liability upfront or spreading the tax payment over thirty years.} This system allows smoothing the wealth transfer over thirty years.

\subsection{Payment in kind}

Liquidity constraint is the objection most often raised against wealth taxes: a taxpayer may be rich on paper yet hold most of their wealth in illiquid form, such as a stake in an unlisted company, real estate or an art collection, that generates little cash, so that a tax payable in cash could force them to sell assets, sometimes at a discount. %and sometimes at the cost of control over their business.

The proposal neutralises it through a graduated set of payment options. First, the administration accepts payment in non-cash liquid assets, in particular shares in listed companies, valued at market price on the day of transfer. A taxpayer whose wealth consists mainly of a quoted equity stake can therefore discharge the tax without selling on the market and without any cash outlay. Second, on a case-by-case basis, the administration may accept payment in illiquid assets, such as shares in non-listed companies, to the extent that the taxpayer's liability exceeds their liquid assets. This mirrors the logic of the Californian bill, which allowed liquidity-constrained taxpayers holding designated highly illiquid assets (typically start-up equity) to attach a claim to those assets rather than face an immediate cash assessment \citep{AB259}.

Payment in shares has an important corollary for the objection that a wealth tax penalises young, capital-intensive firms relative to established ones. % (Section~\ref{sec:objections}).
Because the founder of a start-up can hand over a fraction of their shares rather than cash, the tax does not force the sale of the company %or the dilution of control through a fire sale,
and it does not require the firm to distribute dividends it cannot afford. In case where the tax could make the taxpayer lose control over their company (by reducing their stake below 50\%), the administration may accept payment in non-voting shares.

\subsection{Deferral Accounts}

Unless they settle their entire liability in the first year, taxpayers must open a Deferral Account (DA): a contract with the State to which they attach assets of sufficient value to cover the tax liability. The DA has two functions. It secures the State's claim: the outstanding liability is recovered when the taxpayer dies or moves their tax residence outside the EU, whichever comes first, and the account is closed only once the tax has been paid in full. It supports the payment-in-kind mechanism, since attached assets can be transferred to the tax authority to discharge the liability. %And it provides the enforcement backbone against emigration: because the liability crystallises on exit from the EU, relocating after the 2027 reference date does not extinguish the tax, which sharply limits the fiscal-flight response that has undermined recurrent national wealth taxes.

A missing payment triggers a penalty equal to inflation plus 10\% per year on the overdue amount. This penalty accrues within the DA and is ultimately recovered when the attached assets are sold or generate dividends, so that non-payment is costly but never forces an immediate ownership transfer. The combination of low annual rates, deferral, payment in kind and asset-backed recovery is what makes the tax simultaneously enforceable, non-distortive and respectful of taxpayers' liquidity constraints and asset control.

\section{Revenue estimates}
\label{sec:revenue}

\subsection{Method and assumptions}

Revenue is estimated by microsimulation on the 2023 wealth distributions of the World Inequality Database (WID).\footnote{For most Member States, we use the tabulation that underlies the \href{https://wid.world/world-wealth-tax-simulator/}{WID Wealth Tax Simulator}, which reports, for a grid of euro thresholds, the total net wealth and the number of adults above each threshold. This dataset is the more adequate source, for two reasons: its grid already includes the statutory cut-offs (EUR~1, 10, 100~million and 1~billion), so each bracket's base can be read off directly; and it smooths the top tail, so it resolves the extreme fortunes on which the upper brackets depend.
The eleven smaller Member States that the simulator omits are completed from the raw WID distributional-national-accounts data, which report the average wealth within each of 127 population brackets, up to the top 0.001\%. Applying the schedule to those bracket averages reproduces the lower brackets well but resolves the very top less finely. Because these states are small (their combined contribution is about EUR~0.6~billion), this has no material effect on the EU total.}
Two successive haircuts are applied to the base throughout: 15\% for asset-price depreciation and a further 15\% for evasion and valuation discounts. %, so that the taxable base is reduced to about 72\% of its measured value.

Table \ref{tab:revenue}'s revenue figures are those of the proposed tax's three marginal bands. Column \textit{a} is the 0.1\% levied on the slice of wealth between EUR~1 and 10~million; column \textit{b} the 0.3\% on the slice between EUR~10 and 100~million; and column \textit{c} the 0.5\% on all wealth above EUR~100~million. %, which is the top band of the schedule and so has no upper bound.
The central scenario is the sum of these three bands, $a+b+c$ (also shown as a share of 2023 GDP). Column \textit{d} reports a separate band that does not belong to the central schedule: a 2\% rate on wealth above EUR~1~billion, used only to illustrate a top-heavy variant that would use 2\% as the marginal rate on billionaire wealth, so that the top-heavy figure is $c+\frac{3}{4}d$. All figures are in EUR~million per year.

\subsection{Results}

The central scenario, which combines the three low brackets, raises about EUR~26 billion per year for the EU, or close to EUR~790 billion over the thirty-year horizon. This represents 4.5\% of EU GDP,
in the range of the 5\%-of-GDP benchmark independently proposed by \citet{PisaniFerryMahfouz2023} to finance the French energy transition, and about twice % comparable to the annual servicing cost of
the EU's common (NextGenerationEU) debt. The revenue broadly tracks the distribution of large fortunes across the Union: Germany, France, Italy and Spain together account for 75\% of total revenue. %, and the five largest economies (adding Italy, Spain and Austria) generate around 80\%.
As an own resource, the proceeds would relieve the GNI-based contributions Member States pay into the EU budget; and should the revenue instead endow a European fund (Section~\ref{sec:fund}), a competitiveness-driven mandate would in any case direct much of its investment towards the Member States that host Europe's largest firms, so that the country pattern of contributions and of investment would broadly coincide. %Germany leads the central scenario on the strength of its broad upper-middle wealth (the EUR~1 and EUR~10 million brackets), while France leads on the upper brackets. The top-heavy variant ($c+\frac{3}{4}d$) raises about EUR~25.4 billion per year, drawn from a few thousand centi-millionaires and the EU's billionaires.

These estimates are deliberately conservative for two reasons. They cut the base by two successive 15\% haircuts, both generous. The 15\% allowance for evasion and valuation discounts sits well above the leakage a well-enforced one-off levy should suffer: \citet{Seim2017} finds that the taxable base contracts by at most 1\% in the short run, and \citet{AdvaniTarrant2021} put the medium-run figure at less than 10\% for a 0.5\% wealth tax. The 15\% allowance for depreciation is also cautious, %: \citet{GrindakerVestad2025} report that unlisted shares were on average carried at book value equal to about 70\% of their market value, so a market-value assessment such as ours would, if anything, raise more than the figures shown.
and compounded by our use of 2023 wealth data.
Indeed, the tax would be assessed on wealth held on 1~January~2027:\footnote{Germany prohibits retroactivity (\textit{echte Rückwirkung}), so the assessment date should not precede announcement.} since asset prices have risen since 2023, the true 2027 base, and hence the revenue, would be larger. %The numbers are also static, abstracting from behavioural responses, which for a one-off tax on a past date should be small (Section~\ref{sec:objections}).
%All figures are reproduced by the script \texttt{wealth\_revenue.R} accompanying this appendix, which reads the WID Wealth Tax Simulator dataset and exports the country-by-country tables to \texttt{wealth\_revenue\_tables.xlsx}.

\newpage

\begin{longtable}{lrrrrr|rr}
\caption{Estimated revenue of a one-off wealth tax in the first year, by EU Member State (EUR million per year) %; the base is cut by two successive 15\% haircuts, one for asset-price depreciation and one for evasion, a factor of about 0.72).
The GDP column reports annual central revenue as a share of GDP.%; over the thirty-year horizon the cumulative central levy is about thirty times as large.
}
\label{tab:revenue}\\
\hline
%  &  & \textit{b} & \textit{c} & Central & Central & \textit{d} & Top \\
Country & \makecell{\textit{a}\\0.1\% on\\ 1--10M} & \makecell{\textit{b}\\0.3\% on\\ 10--100M} & \makecell{\textit{c}\\0.5\%\\ $>$100M} & \makecell{Central\\($a+b+c$)} & \makecell{Central\\\% GDP} & \makecell{\textit{d}\\2\%\\ $>$1~bn} & \makecell{Top\\($c+\frac{3}{4}d$)} \\
\hline
\endfirsthead
\multicolumn{8}{c}{\tablename\ \thetable{} -- continued}\\
\hline
 & \textit{a} & \textit{b} & \textit{c} & Central & Central & \textit{d} & Top \\
Country & \makecell{0.1\% on\\ 1--10M} & \makecell{0.3\% on\\ 10--100M} & \makecell{0.5\%\\ $>$100M} & ($a+b+c$) & \% GDP & \makecell{2\%\\ $>$1~bn} & ($c+\frac{3}{4}d$) \\
\hline
\endhead
Germany       & 3{,}050 & 2{,}702 & 2{,}592 & 8{,}344 & 0.20\% & 4{,}963 & 6{,}315 \\
France        & 1{,}994 & 1{,}378 & 3{,}298 & 6{,}670 & 0.24\% & 8{,}812 & 9{,}907 \\
Italy         & 1{,}177 & 448 & 994 & 2{,}618 & 0.13\% & 1{,}327 & 1{,}989 \\
Spain         & 851 & 434 & 866 & 2{,}151 & 0.14\% & 1{,}409 & 1{,}923 \\
Austria       & 296 & 287 & 337 & 920 & 0.19\% & 841 & 967 \\
Netherlands   & 585 & 109 & 220 & 913 & 0.09\% & 182 & 356 \\
Belgium       & 500 & 37 & 186 & 722 & 0.12\% & 527 & 581 \\
Sweden        & 260 & 140 & 312 & 712 & 0.14\% & 829 & 934 \\
Poland        & 169 & 196 & 196 & 560 & 0.08\% & 170 & 323 \\
Denmark       & 199 & 145 & 188 & 533 & 0.14\% & 441 & 519 \\
Czechia       & 74 & 44 & 230 & 348 & 0.11\% & 489 & 597 \\
Ireland       & 247 & 14 & 70 & 331 & 0.06\% & 69 & 122 \\
Portugal      & 187 & 109 & 35 & 330 & 0.12\% & 48 & 70 \\
Luxembourg    & 56 & 37 & 115 & 208 & 0.25\% & 305 & 343 \\
Hungary       & 53 & 62 & 93 & 208 & 0.10\% & 8 & 99 \\
Finland       & 58 & 12 & 82 & 152 & 0.05\% & 84 & 145 \\
Romania       & 39 & 37 & 46 & 121 & 0.03\% & 20 & 61 \\
Greece        & 64 & 17 & 36 & 118 & 0.05\% & 0 & 36 \\
Slovenia      & 18 & 18 & 38 & 75 & 0.11\% & 0 & 38 \\
Cyprus        & 11 & 2 & 33 & 46 & 0.14\% & 0 & 33 \\
Bulgaria      & 13 & 14 & 13 & 40 & 0.04\% & 0 & 13 \\
Croatia       & 18 & 12 & 6 & 36 & 0.05\% & 0 & 6 \\
Lithuania     & 14 & 10 & 7 & 31 & 0.04\% & 0 & 7 \\
Estonia       & 14 & 9 & 6 & 29 & 0.08\% & 0 & 6 \\
Latvia        & 9 & 8 & 3 & 20 & 0.05\% & 0 & 3 \\
Slovakia      & 12 & 7 & 0 & 19 & 0.02\% & 0 & 0 \\
Malta         & 9 & 0 & 0 & 9 & 0.04\% & 0 & 0 \\
\hline
\textbf{EU-27 total} & \textbf{9{,}978} & \textbf{6{,}287} & \textbf{10{,}000} & \textbf{26{,}264} & \textbf{0.15\%} & \textbf{20{,}525} & \textbf{25{,}394} \\
\hline
\end{longtable}

\section{Why this wealth tax can fly}
\label{sec:feasibility}

Having shown that a one-off wealth tax, correctly designed, raises interesting amounts, we now set out why we believe it can actually be adopted. The EU level is more suited than the national one; it does not interact with the taxes Member States already levy; it raises real revenue; %for national budgets; 
it is constitutional in every EU jurisdiction; it rests on solid economic and democratic foundations; and it is increasingly on the political agenda.

\subsection{Suitability of the EU level}\label{sec:level}

A one-off wealth tax is more efficient levied by the Union than by Member States acting separately. Two arguments support this. First, EU law makes it hard for a Member State to secure its tax base against a taxpayer who leaves for another Member State, while the single market makes leaving nearly costless. Second, it is ill-suited to treat a billionaire's fortune as belonging to the country where its owner happens to reside, since that fortune is generated by the worldwide activities of the firms they own.

The instrument a State would use to secure a wealth-tax claim against emigration is an exit tax. Between two Member States, that instrument is tightly constrained by the fundamental freedoms, and the case law has narrowed it steadily. In \emph{de Lasteyrie du Saillant} \citep{CJEULasteyrie2004} the EU Court of Justice struck down the French exit tax on latent gains as a restriction on the freedom of establishment, and held disproportionate the conditioning of deferral on the provision of guarantees or the appointment of a tax representative; the \cite{CJEUN2006} later confirmed that deferral must be granted without security. In \emph{Commission~vs.~Portugal} \citep{CJEUPortugal2016} the Court applied the same reasoning to individuals transferring their residence, finding an unjustified restriction on the free movement of persons and on the freedom of establishment. %\emph{W\"achtler} \citep{CJEUWachtler2019} went further: spreading the charge over several annual instalments was held incapable of removing the cash-flow disadvantage of paying on departure, so that only deferral until the assets are actually disposed of satisfies the proportionality test. 
Germany, which in 2022 replaced its open-ended interest-free deferral for departures within the EU and EEA by a seven-year instalment scheme, is widely thought to have fallen on the wrong side of that line, and the question is now before the Court on a Polish reference \citep{CJEUGena2025}. %Secondary law offers no help: the exit-taxation rule of the Anti-Tax-Avoidance Directive (Article~5 of Directive~(EU)~2016/1164) applies only to taxpayers subject to corporate tax, leaving the exit taxation of individuals entirely to this negative integration \citep{ATAD2016}.

% Two further limits compound the difficulty. Exit taxes reach only the unrealised capital gains embedded in substantial shareholdings, never the wealth stock itself, so even an unimpeded exit tax would give a Member State no way of securing a claim assessed on net wealth. And once 
Once the taxpayer has gone, collection depends on the assistance of the new State of residence under the mutual recovery directive \citep{Directive2010_24}, whose enforcement record is modest.

Departure, meanwhile, is nearly frictionless. Free movement of persons and of capital means that residence and portfolios can be shifted between Member States without exchange controls or administrative hurdles, and several Member States actively bid for wealthy newcomers: Italy offers new residents a substitutive flat charge on all foreign-source income (raised from EUR~100{,}000 to EUR~200{,}000 a year in 2024) and Greece a comparable EUR~100{,}000 lump sum for up to fifteen years. A Member State levying our tax alone would face both this competition and the legal challenges of charging those who respond to it.

Levying the tax at Union level dissolves the problem. % rather than solving it. 
Moving from one Member State to another changes nothing: %, since the liability is fixed by residence on 1~January~2027 and is owed to the Union wherever the taxpayer lives within it, so none of the case law above is even engaged: 
there is no exit, and therefore no restriction to justify. Only departure from the Union crystallises the Deferral Account (Section~\ref{sec:objections}), and there the constraints are far weaker: the freedom of establishment and the free movement of persons do not extend to third countries. %, other than through the EEA Agreement and the free-movement agreement with Switzerland at issue in \emph{W\"achtler}; and what crystallises is in any case a debt already fixed on a past date, not a tax on an unrealised gain. TODO: important
The border that has to be policed is the Union's, which far fewer people cross, and which the Union itself is competent to police.

The second argument is one of principle. The fortunes this levy reaches consist overwhelmingly of stakes in companies whose value is produced across the single market and beyond it. Those firms owe their profits and their valuations to the size of the markets they serve, to free capital movement and to the cross-border enforcement of property rights, all of them Union-wide public goods rather than national ones \citep{Zucman2024}. Attributing the whole of such a fortune to the tax jurisdiction where its owner sleeps is therefore arbitrary, and doubly so once holding structures are taken into account: it hands windfalls to the small Member States that host them and nothing to the countries where the value was actually created. If the base is European in origin, the natural claimant is the European budget.

\subsection{Additionality from existing taxes}
\label{sec:interaction}

The proposed levy would be \emph{additional} to, not a substitute for, the taxes Member States already apply to wealth. Within the EU only Spain still operates a comprehensive recurrent net-wealth tax, and the few other national instruments are narrow. Crucially, all of them share a feature that the present proposal deliberately avoids: they largely spare the business assets that make up the bulk of very large fortunes, so the overlap with our levy is small. % Spain wealth tax revenue: 2G

Spain is the leading case. Its recurrent wealth tax is regionally administered, and several regions (notably Madrid and Andalusia) had granted 100\% relief until the central government introduced, in 2022, a Solidarity Tax on Large Fortunes (ITSGF) to reclaim that base, with marginal rates of 1.7\% above EUR~3~million up to 3.5\% above 11~million. % 1.7% on 3-5M, 2.1% 5-11M, 3.5% > 11M
Yet the Spanish tax exempts business assets: shares held by owner-managers who hold a qualifying stake, play a managerial role and draw most of their income from the firm are relieved, and the primary residence is partly exempt.
On top of this, a cap limits the combined income-and-wealth-tax bill to 60\% of taxable income; but the resulting relief may not exceed 80\% of the wealth-tax liability, so that at least 20\% of it is always due (Article~31 of Law~19/1991). In practice this floor caps the effective wealth-tax rate of the asset-rich but income-poor at one fifth of the statutory rate, about 0.7\% rather than the headline 3.5\%. Together these carve-outs and caps let the very wealthiest largely escape the tax, and avoidance responses to it are well documented \citep{DGTaxud2026}.

France abolished its general wealth tax (ISF) in 2018 and replaced it with a narrow tax on real-estate wealth (IFI), at 0.5\%--1.5\% on net property above EUR~1.3~million, leaving financial and business wealth untaxed.

The Netherlands has no more net-wealth tax: its ``Box~3'' regime taxed a notional return on savings and investments within the income tax, an approach the Supreme Court struck down in 2021 and again in 2024 (Section~\ref{sec:constitutionality}). From 2028 it is to be replaced by the Box~3 Actual Return Act, which taxes the \emph{real} return on wealth: an annual mark-to-market (accrual) tax on the change in value of most assets and debts, combined with a realisation-based capital-gains tax on real estate and on shares in start-ups. This remains a tax on the return to wealth, not on the wealth stock itself, and so leaves the base our levy targets untouched.

Belgium levies a 0.15\% annual tax on securities accounts worth at least EUR~1~million, a narrow charge that reaches listed portfolios but not directly held business stakes.

The proposed levy is therefore complementary rather than duplicative. Because it is one-off, values assets at market and grants \emph{no} business-asset exemption, it falls, for very-high-net-worth individuals, precisely on the unlisted-company and financial holdings that the national taxes spare, while adding only 0.1\% per year at the bottom of the schedule (above EUR~1~million). Where a taxpayer is simultaneously liable to a national recurrent wealth tax (the Spanish solidarity tax or the French IFI), %Member States may choose to credit the two against each other; the baseline
the proposal treats them as cumulative, since even in the most heavily taxed cases the combined annual burden stays below the return on capital and therefore within the bounds courts have treated as non-confiscatory (Section~\ref{sec:constitutionality}).

\subsection{Revenue raising} % for national budgets}

% The revenue is not symbolic. 
At about EUR~26~billion a year (Section~\ref{sec:revenue}), the central scenario raises as much as, or more than, the old national wealth taxes did, despite markedly lower rates, because it corrects their design defects (Section~\ref{sec:objections}). As a genuine own resource, it would reduce the GNI-based contributions Member States pay into the EU budget. % euro for euro, so that the money returns to national treasuries rather than being lost to them. 
This is the property the Council is looking for in a new own resource: a base that is concentrated and additional to what Member States already tax.

\subsection{Constitutionality}
\label{sec:constitutionality}

Wealth taxation is not unconstitutional as such in any EU jurisdiction. The concern that can make a wealth tax fall foul of a constitution is that it becomes confiscatory, taking so large a share of asset returns that it amounts to an expropriation rather than a tax. Low rates dispose of that concern: an annual charge of at most 0.5\% of wealth, only a fraction of the normal return on capital, leaves the principal intact and stays far from any confiscation threshold. The history of national wealth taxes, reviewed below, confirms that they were struck down or abandoned for specific, avoidable defects, unequal valuation, taxation of a fictitious return, or genuinely confiscatory rates, and never because taxing wealth is inherently impermissible.

\paragraph{Germany.} The Federal Constitutional Court suspended the wealth tax in its 1995 decision (2~BvL~37/91) not because wealth taxation is inherently unconstitutional, but because real estate was valued far below market while financial assets were valued at market, violating the equality principle of the Basic Law \citep{BVerfG1995}. The same decision advanced the \emph{Halbteilungsgrundsatz}: the legal principle that wealth and income tax combined should not take away more than roughly half of the potential return of the taxed assets. Two points defuse this. First, the Court itself, in a later ruling, held the \emph{Halbteilungsgrundsatz} not to be legally binding \citep{DGTaxud2026}. Second, the proposed tax would satisfy it in any case: an annual charge of at most 0.5\% of wealth is only a small fraction, of the order of a tenth, of the normal return on capital, so that even added to income tax the combined burden stays far below half of that return. A uniformly market-valued, low-rate one-off tax thus avoids both defects that felled the German tax.

\paragraph{Spain.} The Constitutional Court \emph{upheld} the central solidarity wealth tax in December 2023 against challenges from several regions, confirming that a wealth tax is compatible with the Spanish constitution \citep{STC2023}.

\paragraph{The Netherlands.} %The Dutch case is often misread.
The Netherlands abolished its genuine net-wealth tax in 2001 for practical, not legal, reasons: the old tax was widely avoided or evaded and raised little. The wealth tax was folded into the income tax as ``Box~3'', which taxes a \emph{notional} (assumed) return rather than actual wealth or actual income, a design chosen for administrative simplicity, not because taxing wealth was thought unconstitutional. It is precisely the taxation of a \emph{fictitious} return that the Supreme Court struck down, in the 2021 ``Christmas Judgment'' and again in 2024, as an infringement of the right to property and the prohibition of discrimination under the European Convention on Human Rights \citep{HogeRaad2021}. The Netherlands is now moving to a tax on \emph{actual} returns from 2028. A charge levied on \emph{actual} wealth, with an ability-to-pay safeguard, does not share this flaw; if anything the Dutch saga shows that taxing real wealth is on firmer legal ground than taxing an imputed one.

\paragraph{France.} The Conseil constitutionnel has repeatedly upheld wealth taxation (the ISF and later the IFI), on the condition that the burden not become confiscatory, which for a holding tax may require an income-based cap only when the rate is high \citep{CPO2018, ConseilConstitutionnel2012}. The threshold is telling: in the ISF case-law a rate of 0.5\% was validated without any cap, whereas a rate of 1.8\% required one \citep{CPO2018}.
The maximum annual rate proposed here is 0.5\%, the very level the Conseil accepted with no cap, so the French standard is comfortably met. It is worth noting that the Conseil constitutionnel also judged that ability to pay may be assessed on wealth ownership rather than income flows (Decision n\textdegree~2019-769~QPC), confirming the legitimacy of a wealth tax.

\paragraph{Austria.} Austria discontinued its wealth tax in 1993 by an act of government, not of any court, and for reasons of design rather than principle: fear of capital flight, widespread evasion, high administrative cost, and low and declining revenue, compounded by the unequal treatment of grossly under-valued real estate (the same defect that later led the Constitutional Court to strike down the Austrian inheritance tax) \citep{DGTaxud2026}.

However, one genuine legal obstacle is specific to Austria: Article~1(1)(2) of the constitutionally entrenched Final Taxation Act (\emph{Endbesteuerungsgesetz}) bars any wealth tax on bank deposits and domestically issued bonds. Overcoming it would require either amending that constitutional law by a two-thirds majority, or granting Member States the option to exempt those asset classes, which in any case make up only a small share of large fortunes. %The Austrian experience is thus a catalogue of avoidable design failures, not evidence that wealth cannot be taxed.

% \paragraph{European law.} Two European-level points complete the picture. First, Article~1 of Protocol No.~1 to the European Convention on Human Rights protects the peaceful enjoyment of possessions but explicitly permits States to levy taxes, subject to a ``fair balance'' between the general interest and the individual burden; the Court finds a violation only where the charge is disproportionate or confiscatory, and a one-off levy of at most 0.5\% of wealth per year with an ability-to-pay ceiling %(the taxpayer never owes more than the schedule applied to current wealth)
% sits well within that margin. Second, %and importantly for the sovereign-wealth-fund destination of the revenue,
% Article~345 of the Treaty on the Functioning of the European Union provides that ``the Treaties shall in no way prejudice the rules in Member States governing the system of property ownership'': EU law is neutral as between private and public ownership, so using the proceeds to build up collective, public ownership of productive capital raises no objection under the Treaties.

The recurring constitutional theme across jurisdictions is that wealth taxes fail when they value assets unequally, tax fictitious returns, or become confiscatory, never simply because they tax wealth. The present design is built to avoid all three failure modes.

\subsection{Solid economic justification}

The instrument rests on solid economic foundations. \citet{Farhi2010} shows that a one-off capital levy is
% in the canonical model of optimal taxation with aggregate shocks and incomplete markets,
the efficient response to a negative fiscal shock, % is a one-off capital levy, exactly the contingency created by the shocks that motivate new own resources:
precisely the case created by
the pandemic-era debt, the security build-up and the financing of the transition. Because the base is a past stock, the levy raises revenue without blunting the incentive to work, save or invest after 2027, the classical efficiency property of a lump-sum tax. Progressive wealth taxation is independently justified on distributive grounds: \citet{SaezZucman2019} argue that %, by reaching the \emph{stock} of fortune rather than the \emph{flow} of realised income,
it is the most appropriate tool for restoring progressivity at the very top, since the ultra-wealthy report little taxable income.

The design also has concrete precedents. It adapts the deferral and one-off mechanism of the Californian wealth-tax bill \citep{AB259}, which let liquidity-constrained taxpayers attach a claim to designated illiquid assets rather than face an immediate cash assessment, and it adopts the low-rate one-off logic of \citet{PisaniFerryMahfouz2023}, who recommended an exceptional levy of 5\% of GDP spread over thirty years to finance the climate transition. These are not speculative constructs but instruments that leading economists have already drafted and defended.

\subsection{Democratic appeal}

Beyond efficiency, the tax answers a democratic demand. Roughly two thirds of European citizens support taxing the very wealthy \citep{TaxObservatory2025}, so an own resource built on large fortunes rests on a broader popular mandate than most alternatives on the Council's table. This support extends to the cross-border dimension that an EU-level resource entails: representative surveys across high-income countries find majority backing for international redistribution \citep{Fabre2025} and for global redistributive and climate policies \citep{FabreDouenneMattauch2025}, so a European wealth levy channelled to common objectives aligns with, rather than runs ahead of, public opinion. It is also a matter of basic fairness: large fortunes are today taxed at effective rates well below those borne by ordinary households \citep{Zucman2024}, and a levy that reaches the stock of wealth restores the progressivity that income taxes alone can no longer deliver at the top. Finally, the extreme concentration of wealth is itself a concern for democracy, since it concentrates economic and political power; a modest, broadly supported levy that returns part of that wealth to the common budget strengthens rather than threatens democratic legitimacy.

\subsection{Suitability with the political debates}

Taxing large fortunes in a coordinated way is no longer a marginal idea. At the request of the Brazilian G20 presidency, \citet{Zucman2024} presented to the July~2024 Rio meeting of G20 finance ministers a blueprint for a coordinated minimum tax on ultra-high-net-worth individuals, under which billionaires would pay at least 2\% of their wealth each year.\footnote{The G20 Rio de Janeiro Ministerial Declaration on International Tax Cooperation (July~2024) committed member states, ``with full respect to tax sovereignty'', to ``seek to engage cooperatively to ensure that ultra-high-net-worth individuals are effectively taxed'', through the exchange of best practices, debate on tax principles and the design of anti-avoidance mechanisms; the commitment was reaffirmed in the leaders' declaration of November~2024.} Several governments, including Brazil, South Africa, France and Spain, have supported the coordinated approach. Spain, the only Member State still operating a net-wealth tax (which raised about EUR~2.2~billion in 2023), has been its most consistent EU advocate.

At European level, the instrument is now under serious institutional consideration. The European Commission has commissioned a study on the effectiveness of wealth-related taxes targeting ultra-high-net-worth individuals across EU and non-EU countries, and the European Parliament's Subcommittee on Tax Matters (FISC) held a public hearing on the taxation of ultra-high-net-worth individuals in December~2025. %The idea that the revenue should feed the Union's own resources is not ours alone: it is exactly the framing of the European Citizens' Initiative ``Taxing great wealth to finance the ecological and social transition'', registered by the Commission in 2023 and %backed by economists including Thomas Piketty and Joseph Stiglitz, 
% which called on the EU to introduce a tax on great wealth contributing to the Union's own resources. 
A one-off EU levy of the kind proposed here offers a concrete, legally robust way to act on this momentum.

\section{Responses to main objections}
\label{sec:objections}

\subsection{Expatriation}

The objection most often raised against wealth taxes is that the wealthy will simply leave. This concern is real for a recurrent tax, %whose base can walk out of the country year after year,
but it has little force against a one-off levy assessed on residence at a fixed past date. Because the liability is fixed by where a taxpayer lived on 1~January~2027, no later move can change what is owed. %: an individual resident in the EU on that date remains liable in full even after relocating.
Emigration therefore cannot erase the tax, only the ease of collecting it, and that is what the Deferral Account is designed to secure. The outstanding balance is attached to assets and crystallises when the taxpayer dies or transfers their tax residence outside the EU, so that leaving triggers, rather than avoids, payment. This is the same logic as an exit tax, applied to a liability that was already fixed before any departure.

% The migration responses documented in the empirical literature are, in addition, modest and specific to recurrent taxes. The often-cited Swiss and Spanish studies \citep{Brulhartetal2022, DuranCabreetal2019} measure reactions to permanent annual taxes, and even there the estimated flows are small relative to the base. The French \citet{CAE2025} note reaches the same conclusion for high-income taxpayers: wealth taxation accelerates fiscal exile only marginally, with no notable macroeconomic effect. Applied to a one-off tax on a past stock, these elasticities overstate the true leakage by a wide margin, because the mechanism they capture, namely adjusting future residence to lower a future recurring bill, does not operate here. The combination of a backward-looking base and an exit-triggered Deferral Account thus removes the channel through which fiscal flight erodes recurrent wealth taxes.

\subsection{Liquidity}

Another common objection is that wealthy taxpayers, though rich on paper, may lack the cash to pay. The design answers this on two levels: the annual rate is low (at most 0.5\% of wealth, payable out of normal capital income rather than principal), and payment can be made in listed shares or, in genuinely illiquid cases, in kind,
so that no market sale or cash outlay is required. % Empirically, liquidity is rarely binding at these rates: a 2\% annual charge, four times the top rate here, is already below the average return on large portfolios \citep{TaxObservatory2025, WorldInequalityReport2022}. The low rates and the payment-in-kind option therefore dissolve the liquidity objection.

% \subsection{Harm to start-ups relative to established firms}

% A related concern is that a wealth tax bears more heavily on the founders of young, cash-poor firms than on the owners of mature, dividend-paying companies, because the former cannot easily generate the cash to pay. The payment-in-shares option is the direct remedy: a founder can discharge the liability by transferring a fraction of equity rather than by extracting cash from a firm that needs to reinvest it. The Deferral Account reinforces this, allowing the liability to sit against the illiquid stake until a liquidity event (a sale, an IPO, or a dividend) occurs \citep{AB259}. The one-off nature of the tax further protects entrepreneurship: only wealth accumulated by 2027 is taxed, so the returns to founding and growing a company after that date are untouched.
% % https://x.com/sc_cath/status/1933169423826698735

\subsection{Reduced investment and production}

A wealth tax could in principle reduce private investment by lowering the after-tax return to capital and by consuming tax revenue through public spending, eventually lowering GDP. This concern does not apply here for two reasons. First, the tax is one-off and backward-looking, so it does not alter the marginal return on any investment decided after 2027. Second, if the revenue is used to endow a European sovereign fund rather than consumed (as proposed in Section~\ref{sec:fund}), assets withdrawn from private portfolios are re-invested, not spent, so the aggregate capital stock is preserved and ownership merely shifts from private to collective hands. On that route the fund's returns are designed to more than cover the servicing of the common debt, turning a one-off levy into a permanent public asset.
% https://x.com/sc_cath/status/1967435361585344722
% https://x.com/LevyAntoine/status/1962943296701186088

\subsection{Deterring entrepreneurs}

The claim that wealth taxes deter entrepreneurship applies to \emph{recurrent} taxes that repeatedly claw back the fruits of success. It does not apply to a one-off tax on a past stock. An entrepreneur contemplating a new venture in 2028 faces no wealth-tax consequence from success, because the reference date has passed. The incentive to create, take risks and build is therefore left intact, which is precisely why \citet{PisaniFerryMahfouz2023} and the classical public-finance literature on capital levies favour exceptional over permanent instruments.

\subsection{Valuation}

A recurring objection is that wealth is hard to value, especially unlisted businesses, real estate and art. The one-off design blunts this objection at the root: because the tax is assessed on a single reference date, assets need to be valued only once, exactly as they already are at death for inheritance tax, rather than re-valued every year. The recurring administrative cost that weighs on permanent wealth taxes therefore does not arise.

The experience of countries that value wealth annually shows the task is feasible but must be done at market value. Norway, which has levied a net-wealth tax without interruption since 1892, values listed shares at market, primary dwellings by hedonic pricing and other real estate by construction cost, and unlisted firms at their book value excluding intangible assets \citep{DGTaxud2026}. That last convention systematically under-values young, intangible-intensive companies, the very firms whose value lies in their brand, technology or growth prospects, so Norway's base is, if anything, too generous to business owners \citep{GrindakerVestad2025}. We therefore assess all assets at fair market value. %Where a market value is genuinely uncertain, the payment-in-kind option turns the difficulty into a safeguard: a taxpayer who claims that their unlisted shares are worth little may settle the tax by handing over those very shares at their self-declared value, which removes the incentive to understate.
% The historical failures of wealth taxation came not from valuing assets at market but from the opposite: Germany and Austria taxed grossly under-valued real estate alongside market-valued financial assets, and it was that inequality of treatment, not valuation as such, that the courts condemned (Section~\ref{sec:constitutionality}).

\subsection{The low yield of past wealth taxes}

Skeptics point out that recurrent wealth taxes have historically raised little, around 0.2\% of GDP in France (ISF) and Spain, and roughly 0.1\% in the former German and Austrian taxes \citep{DGTaxud2026}. This is a critique of how those taxes were designed, not of wealth taxation as such. Their yield was hollowed out by three features that the present proposal rejects. First, generous exemptions for business assets. %: France (through the Dutreil pact and the exemption of professional assets), Spain (shares of qualifying owner-managers) and Sweden (large closely-held firms) all spared precisely the assets that dominate the largest fortunes, so the tax fell on the upper-middle class while the very rich escaped.
Second, income-based caps: for instance, France's cap let the billionaires %wealthiest 0.001\%
be taxed at an effective rate of less than 0.1\% \citep{DGTaxud2026}. %, and Spain caps the bill at 60\% of income \citep{DGTaxud2026}.
Third, annual assessment on a mobile base, which invited migration and avoidance year after year. Our levy grants no business-asset exemption, has no income cap (it is instead capped by an ability-to-pay ceiling on \emph{current wealth}), and is assessed once on a past stock that cannot emigrate. It is precisely because it corrects these three defects that it raises as much revenue as old taxes did despite significantly lower tax rates. % multiples of what the old taxes did---about EUR~24.5~billion a year for the central scenario, against the roughly EUR~3~billion Spain's tax yields today.

% \subsection{Fiscal flight and the Levy--Catherine critique}

% The sharpest recent objections come from Antoine Levy (UC Berkeley) and Sylvain Catherine (Wharton), who, together with a group of economists, contested the revenue estimates of the ``Zucman tax'' in France. Their arguments are of three kinds. First, that a large share of the apparent tax base consists of retained earnings inside holding companies that have already borne corporate income tax, so that a wealth tax amounts to double taxation and its true net yield is far smaller than advertised---on their reading, closer to a few billion than to EUR~20 billion for France. Second, and most prominently in Levy's public thread, that behavioural responses---above all fiscal exile of the very wealthy---would turn the measure into a \emph{net loss} for public finances, which he provocatively quantifies at EUR~30 billion \citep{LevyThread2025}. Third, that the valuation of unlisted and illiquid assets is arbitrary and open to manipulation.

% Our design answers each point. On \emph{double taxation}, the objection conflates a tax on income flows with a tax on a wealth stock; a one-off levy on the stock held in 2027 is levied once, at a low rate, and does not recur on the same retained earnings year after year. On \emph{fiscal flight}, the Levy critique is fundamentally a critique of a \emph{permanent} tax, whose base can walk out of the country every year. A one-off tax assessed on residence at a fixed past date cannot be escaped by subsequent emigration, and the Deferral Account crystallises the liability on exit from the EU, so the migration elasticities that drive Levy's ``net loss'' figure do not bite. The empirical migration responses documented for recurrent taxes \citep{Brulhartetal2022, DuranCabreetal2019} therefore overstate the leakage of a one-off levy by a wide margin. It is also worth noting that the debate is far from settled among economists: the note of the \citet{CAE2025} finds that wealth taxation accelerates the fiscal exile of top earners only marginally, with no notable macroeconomic consequence, and Olivier Blanchard, among others, has publicly defended the feasibility of taxing the very wealthy. On \emph{valuation}, the payment-in-kind option turns the difficulty into a safeguard: a taxpayer who claims their unlisted shares are worth little may settle the tax by handing over those very shares at their self-declared value, which sharply curbs the incentive to understate.

% NOTE: inflation in table
% NOTE: check and answer Andrews et al. 26

\section{A more strategic use: a European sovereign fund}
\label{sec:fund}

Everything above establishes that the levy works as an own resource: its proceeds could be paid into the EU budget, relieving Member States' contributions, or used to repay the common NGEU debt. But we think there is an even better way to deal with the debt-reimbursement problem. Rather than shrink the Union's balance sheet by repaying the debt, the proceeds could be invested, at a time when major investments are needed for the ecological transition and for competitiveness.

The idea is to endow a European sovereign wealth fund. Such funds are no longer the preserve of resource-rich states: China, Singapore, Australia and South Korea each run one worth more than 10\% of their GDP, %. Some of the largest sovereign funds in the world are run by economies with little or no natural-resource windfall: China (14\% of GDP across its state investment vehicles and four times its GDP in central or local government-owned enterprises), Singapore (about 240\% of its GDP), Australia (13\% of GDP), and South Korea (12\% of GDP),
% In China, 74% of private equity fundraising comes from government sources https://x.com/NikkeiAsia/status/2081265787881324698/photo/1
and Canada launched a national fund in 2026 financed by public debt. %Norway's fund, the largest of all, exceeds four times its GDP.
Worldwide, sovereign wealth funds manage assets worth close to 13\% of global GDP, yet the European Union has no comparable fund. Sized at 5\% of EU GDP, the levy would build up a fund whose assets would eventually reach about twice the common debt, so that its long-run dividends would exceed the interest owed several times over. As the fund reinvests the proceeds rather than consuming them, the aggregate capital stock is preserved and ownership merely shifts from private to collective hands (Section~\ref{sec:objections}).

\subsection{Governance and mandate of the fund}

The proceeds would endow a European sovereign wealth fund under democratic control. Its board would be appointed by the political groups of the European Parliament, each group naming board members and exercising a voting right proportional to its number of seats in the Parliament, so that the fund's direction reflects the Union's democratic balance rather than the discretion of a single executive.

The fund's mandate would be threefold: to hold a broadly diversified portfolio of equities, so as to spread risk and secure a stable return; to steer the ecological transition; and to select strategic investments that strengthen Europe's future-preparedness. On the model of the recent Canadian national fund, the mandate would give
priority to \emph{primary} investments (initial public offerings, capital increases and other subscriptions of newly issued equity) in European companies with a strategic dimension or a sustainable project, so that the fund channels fresh capital into productive investment rather than mostly reallocating ownership of existing shares.

\section{Possible variants}
\label{sec:variants}

Several variants are compatible with the architecture above but are left for future work and not developed here: (i) converting the one-off levy into a \emph{permanent} recurrent wealth tax; (ii) applying \emph{higher rates} above EUR~100~million wealth, closer %to the top-heavy scenario of Table~\ref{tab:revenue} or
to the 2\%--3\% minimum-tax proposals of \citet{Zucman2024}; and (iii) pairing the wealth tax with a reform of \emph{inheritance} and \emph{capital gains taxation}, which reaches the same tax base at the point of transmission. %; and (iv) dedicating the revenue to repaying the common debt or to the EU budget rather than to a sovereign wealth fund. 
Each of these would change the efficiency, legal and political properties of the instrument and deserves separate analysis.
% NOTE
% taxer à la place les plus-values latentes (plan Biden)
% coopération renforcée
% - la simple hausse de l'ISF norvégien de 0.85% à 1.1% a suscité tant de départ que les recettes fiscales ont diminué- Une exit-tax assise sur les plus-values latentes est en effet contournable par anticipation, via des restructurations préventives ou des transferts progressifs en holding, par une population qui dispose du conseil juridique et du temps nécessaire pour les organiser- Selon les travaux récents de Rexecode (Didier et Koléda, Document de travail n°99, janvier 2026), les recettes fiscales perdues par érosion de la base de l’ISF auraient atteint environ 9 milliards d’euros par an, contre 2 à 5 milliards encaissés - soit un rendement net négatif.

% - rendre le taux à 10% un deferral penalty autorisé en cas d'illiquidité et le 3/5/7 un normal deferral charge

% \small
\bibliographystyle{apalike}
\bibliography{wealth_references}

\end{document}


% 1. Faible efficacité fiscale: The Economist affirme qu'il est très difficile d'évaluer la valeur réelle des actifs (biens immobiliers, œuvres d'art, participations commerciales), ce qui rend la taxe coûteuse à gérer et peu rentable. Mais est-ce vrai? Ce n'est pas si difficile d'évaluer la valeur réelle d'actifs comme ceux illustrés entre parenthèses, n'est-ce pas? 

% 2. Faibles recettes: The Economist affirme que les pays qui ont expérimenté cette taxe ont constaté qu'elle générait un faible rendement par rapport à son coût de mise en œuvre. Mais a-t-elle vraiment un coût opérationnel plus élevé? Est-elle vraiment source de bureaucratie supplémentaire??

% 3. Fuite des capitaux: The Economist utilise ici l’excuse classique selon laquelle les très riches changent facilement de pays ou transfèrent des actifs, réduisant encore les recettes fiscales. Mais est-ce vraiment le cas, statistiquement parlant? 

% 4. D’autres impôts sont déjà progressifs: The Economist affirme que des impôts progressifs existent déjà dans le monde entier, via l’impôt sur le revenu, l’impôt sur les plus-values ​​et les droits de succession, qui font déjà payer proportionnellement plus cher les plus riches. Mais alors, quelle est la différence avec l’impôt sur les super-riches ? Devrions-nous, au Brésil, d’abord augmenter la progressivité de notre impôt sur le revenu au lieu d’un impôt sur les super-riches ? Et avons-nous réellement des droits de succession et des droits de succession ici?

% 5. Effet symbolique et populiste: The Economist affirme que les impôts ont plus de valeur morale que pratique. Ils rapportent peu et peuvent être source de frustration lorsqu’ils n’améliorent pas les services publics. Je pensais déjà que c’était absurde kkkk mais n’hésite pas à commenter.

% 1. L'évaluation des actifs est faite au moment des successions. Elle est aussi faite chaque année en Suisse et en Norvège, où il y a une taxe sur la fortune. C'est donc faisable.
% 2. Complètement faux. C'est dû au fait que les pays qui avaient un impôt sur la fortune (comme la France) le faisait mal, avec plein d'exemptions, ce qui résultait dans un évitement de l'impôt de la part des plus riches, les milliardaires. La proposition de Zucman apprend des erreurs du passé et évite ces écueils.
% 3. Là encore, le système fiscal peut être conçu pour éviter l'expatriation, grâce à une exit tax et en rendant les contribuables redevables de la taxe plusieurs années après leur départ.
% 4. Je sais pas s'il y a des droits de succession au Brésil mais il faut augmenter la progressivité des trois : revenu (y.c. plus-values), fortune, succession. Les milliardaires n'ont pas de revenu donc le revenu ne suffit pas. Dans une certaine mesure les successions peuvent se substituer à la fortune, mais pourquoi pas faire les deux (ne pas attendre que les riches meurent pour les taxer) ? 
% 5. Absurde. Ça rapporte de l'argent et peut donc améliorer les services publics.
% - ça va réduire l'investissement (à moins que les recettes soient réinvesties dans des actions, auquel cas elle ne rapporte rien au budget, mais permet à l'État de prendre un peu contrôle sur les entreprises) => certes mais on a probablement plus besoin d'investissement dans l'éducation et la transition écologique que dans les grosses entreprises qui ont déjà plein d'argent.
% - ça va forcer les propriétaires de sociétés non cotées à vendre des parts pour pouvoir payer l'impôt, or c'est pas liquide, ils vont pas pouvoir et vont perdre le contrôle de leur entreprise. => l'État peut accepter le paiement en actions plutôt qu'en cash. Certes ils perdront le contrôle de leur entreprise s'ils font ça, mais c'est le but. Veut-on que l'IA soit contrôlée par une personne ou par le public ?
% - c'est pas constitutionnel (du moins en France) => possible, on sait pas trop avant que le Conseil Constitutionnel ne statue. Mais c'est pas un argument contre la taxe, c'est un argument pour changer la Constitution ou avoir un mouvement social si important qu'il fasse bouger les lignes.
